\section{Related Work}

We briefly discuss related work on differentiable rendering, followed by a broad overview of embodied visual navigation, and finally, we focus on the work most relevant to us: Instance ImageGoal Navigation (IIN).

\textbf{Differentiable Rendering.} To achieve photo-realistic scene capture, differentiable volumetric rendering has gained prominence with the introduction of Neural Radiance Fields (NeRF) \cite{mildenhall2020nerf}. NeRF utilizes a single Multi-Layer Perceptron (MLP) to represent a scene, performing volume rendering by marching along pixel rays and querying the MLP for opacity and color. Due to the inherent differentiability of volume rendering, the MLP representation is optimized to minimize rendering loss using multi-view information, resulting in high-quality novel view synthesis (NVS). The primary limitation of NeRF is its slow training speed. Recent advancements have addressed this issue by incorporating explicit volume structures, such as multi-resolution voxel grids \cite{fridovich2022plenoxels, liu2020neural, sun2022direct} and hash functions \cite{muller2022instant}, to enhance performance.


In contrast to NeRF, 3DGS \cite{kerbl3Dgaussians} employs differentiable rasterization. Unlike ray marching, which iterates along pixel rays, 3DGS iterates over the primitives to be rasterized, similar to conventional graphics rasterization. By leveraging the natural sparsity of a 3D scene, 3DGS achieves an expressive representation capable of capturing high-fidelity 3D scenes while offering significantly faster rendering. 
Comprehensive review of the developments in 3D Gaussian Splatting \cite{chen2024survey} highlights the method's versatility and applications across various domains.
Leveraging these advantages, a growing body of research begins to explore various innovations, including deformable or dynamic Gaussians \cite{luiten2023dynamic, wu20234dgaussians, yang2023deformable3dgs}, advancements in mesh extraction and physics simulation \cite{feng2024gaussian,xie2023physgaussian,guedon2023sugar}, as well as applications in Simultaneous Localization and Mapping (SLAM) \cite{keetha2023splatam, matsuki2023gaussian, yugay2023gaussianslam}. This surge of interest in the capabilities of 3DGS leads us to consider its potential for embodied visual navigation. Given that the Gaussian representation inherently encapsulates explicit scene geometry and the parameters required for rendering, we posit that 3DGS can enhance decision-making in map-based visual navigation. The explicit nature of the Gaussian representation provides a rich, condensed form of environmental data, which can be effectively utilized to inform and guide autonomous agents in navigation.

\textbf{Embodied Visual Navigation.} Embodied Visual navigation includes several topics: ObjectGoal Navigation (ObjectNav), Multi ObjectGoal Navigation (MultiON), ImageGoal Navigaiton (ImageNav), and Instance ImageGoal Navigation (IIN). ObjectNav \cite{majumdar2022zson, chaplot2020object, wang2024find, chang2020semantic, chaplot2021seal} requires an agent to navigate to any instance of a specified object category within the environment. MultiON \cite{chen2022learning, wani2020multion, schmalstieg2022learning, marza2023multi}, on the other hand, tasks the agent to sequentially navigate to a series of objects. ImageNav \cite{al2022zero, choi2021image, yadav2023ovrl, yadav2023offline, Kwon_2023_CVPR,Chaplot_2020_CVPR,savinov2018semi,TSGM, wasserman2022lastmile} involves navigating to the camera pose from which a target image is captured. In contrast, IIN \cite{krantz2023navigating, bono2023endtoend, lei2024instanceaware} requires navigating to the specific instance captured by the camera in the target image. Collectively, these navigation tasks span a spectrum from semantic-level navigation in ObjectNav and MultiON to fine-grained instance-level navigation in ImageNav and IIN, comprehensively capturing the problem space of embodied visual navigation. 

Numerous approaches to solving embodied visual navigation utilize deep reinforcement learning (DRL) to develop end-to-end policies that map egocentric vision to action \cite{al2022zero, majumdar2022zson, yadav2023offline, zhu2017target}. However, acquiring skills related to visual scene understanding, semantic exploration, and long-term memory are challenging in an end-to-end framework. Consequently, these methods often incorporate a combination of careful reward shaping \cite{choi2021image}, pre-training routines \cite{yadav2023offline}, and advanced memory modules \cite{savinov2018semi, mezghan2022memory, wu2019bayesian}. In contrast to end-to-end DRL, alternative approaches decompose the problem into sub-tasks that can be optimized in a supervised manner. These sub-tasks include graph prediction via topological SLAM \cite{Chaplot_2020_CVPR}, graph-based distance learning \cite{hahn2021no, shah2021rapid}, and camera pose estimation for last-mile navigation \cite{wasserman2022lastmile}. Chaplot \etal \cite{chaplot2020object} decompose the embodied visual navigation task into exploration, object detection, and local navigation. Building on this, CLIP on Wheels (CoW) \cite{gadre2023cows} utilizes a similar decomposition strategy, focusing on exploration and object localization to handle an open-set object vocabulary. Modular approaches also show promise for effective simulation-to-reality transfer (Sim2Real). Gervet \etal \cite{gervet2023navigating} conduct a Sim2Real transfer of both modular and end-to-end systems, demonstrating that modularity effectively mitigates the visual Sim2Real gap that impairs the performance of end-to-end policies. 

Modular approaches typically represent the environment using a map and utilize this map to acquire the knowledge necessary for navigation within that environment. Bird's-Eye View (BEV) maps~\cite{liu2023bird, chaplot2020object, krantz2023navigating, lei2024instanceaware} are a form of metric map that project the entire scene from an overhead perspective, representing the occupancy status, exploration state, and semantic categories of specific areas. This map representation encodes information about each region into a grid, achieving accurate spatial awareness. In contrast, topological maps~\cite{Chaplot_2020_CVPR, TSGM, Kwon_2023_CVPR} focus more on describing the spatial relationships between different areas in a scene. This approach allows agents to navigate based on the connectivity of spaces rather than relying solely on geometric coordinates. To represent a scene more meticulously, 3D-aware maps have been proposed~\cite{chaplot2021seal, zhang20233d, tan2024self, liu2024volumetric}. Incorporating 3D consistency can enhance perception in navigation. Building upon 3D maps, we further advance by adopting 3D Gaussian Splatting (3DGS)~\cite{kerbl3Dgaussians} as a novel map representation. This enables the map to synthesize appearance views of specific object instances, which we have demonstrated to be effective in the Instance ImageGoal Navigation task.


\textbf{Instance ImageGoal Navigation.} The IIN task introduces distinct challenges compared to the ImageGoal Navigation. First, in IIN, the goal image must depict a specific object instance. In contrast, ImageGoal Navigation may use randomly captured photos that could include insignificant elements, such as large white walls. Second, the camera parameters used to capture the goal image do not necessarily match those of the agent's camera. Therefore, to succeed in IIN, an agent must be adept at identifying the target object among numerous candidates of the same class with goal object and recognizing it from various viewpoints. To address the above challenges, Krantz \etal \cite{krantz2023navigating} develop a general pipeline for aligning the same object from different angle of views. Bono \etal \cite{bono2023endtoend} present an end-to-end approach in the IIN task, while Lei \etal \cite{lei2024instanceaware} propose a method that mimics human behaviour for verifying objects at a distance. Existing methods focus partly on designing sophisticated modules \cite{krantz2023navigating, lei2024instanceaware} and partly on pre-training on pretext tasks \cite{bono2023endtoend}. Our approach differs in that we concentrate on designing a new map representation. Through this novel form of map, we can better establish the connection between target description and target locations, thereby facilitating visual navigation.


